\documentclass{article}

\usepackage{amsmath,amssymb}
\usepackage{xurl}
\usepackage[colorlinks=true,linkcolor=blue,citecolor=blue,urlcolor=blue]{hyperref}

% Shared commands for notes in docs/paper-gaps/.

\DeclareUnicodeCharacter{2080}{\ifmmode{}_{0}\else\textsubscript{0}\fi}
\DeclareUnicodeCharacter{2081}{\ifmmode{}_{1}\else\textsubscript{1}\fi}
\DeclareUnicodeCharacter{2082}{\ifmmode{}_{2}\else\textsubscript{2}\fi}
\DeclareUnicodeCharacter{2099}{\ifmmode{}_{n}\else\textsubscript{n}\fi}

\newcommand{\Lean}{\textsc{Lean}}
\ExplSyntaxOn
\NewDocumentCommand{\leanid}{m}{
  \ifmmode
    \textsf{\detokenize{#1}}
  \else
    \group_begin:
    \urlstyle{sf}
    \tl_set:Nn \l_tmpa_tl {#1}
    \tl_replace_all:Nnn \l_tmpa_tl {\_} {_}
    \exp_args:NV \path \l_tmpa_tl
    \group_end:
  \fi
}
\ExplSyntaxOff
\providecommand{\tr}{\operatorname{tr}}
\newcommand{\tnleanIssueBase}{https://github.com/LionSR/TNLean/issues/}
\newcommand{\tnleanPrBase}{https://github.com/LionSR/TNLean/pull/}
\newcommand{\tnleanDocBase}{https://sirui-lu.com/TNLean/docs/}
\newcommand{\ghissue}[1]{\href{\tnleanIssueBase#1}{GitHub issue \##1}}
\newcommand{\ghpr}[1]{\href{\tnleanPrBase#1}{PR \##1}}
\newcommand{\doclink}[1]{\href{\tnleanDocBase#1.html}{\path{#1}}}

% Dirac notation and the identity operator, as in blueprint/src/macros/common.tex.
% \providecommand keeps note-local definitions authoritative.
\usepackage{dsfont}
\providecommand{\ket}[1]{|#1\rangle}
\providecommand{\bra}[1]{\langle #1|}
\providecommand{\braket}[2]{\langle #1 | #2 \rangle}
\providecommand{\ketbra}[2]{|#1\rangle\!\langle #2|}
\providecommand{\Id}{\mathds{1}}
\providecommand{\one}{\mathds{1}}
\providecommand{\C}{\mathbb{C}}
\providecommand{\MN}[1]{M_{#1}(\C)}
\providecommand{\MD}{M_D(\C)}

% External referencing for paper sources.  The build helper writes sanitized
% auxiliary files under build/paper-aux/ and excludes duplicated source labels.
\usepackage{xr}

% Import a generated paper auxiliary file with a caller-supplied label prefix.
% The candidate paths support builds from docs/paper-gaps/, the repository root,
% and build/paper-aux/ itself.
\newcommand{\importpaperaux}[2]{%
  \IfFileExists{../../build/paper-aux/#2.aux}{%
    \externaldocument[#1]{../../build/paper-aux/#2}%
  }{%
    \IfFileExists{build/paper-aux/#2.aux}{%
      \externaldocument[#1]{build/paper-aux/#2}%
    }{%
      \IfFileExists{#2.aux}{%
        \externaldocument[#1]{#2}%
      }{}%
    }%
  }%
}

% General external reference macros
\newcommand{\paperref}[2]{\ref{#1-#2}}
\newcommand{\papereqref}[2]{(\ref{#1-#2})}

% CPSV16 source (1606.00608 / MPDO-22-12-17-2.tex) external document and shortcuts
\importpaperaux{cpsv16-}{1606.00608/MPDO-22-12-17-2-tnlean}
\newcommand{\cpsvref}[1]{\paperref{cpsv16}{#1}}
\newcommand{\cpsveqref}[1]{\papereqref{cpsv16}{#1}}

% Verdict marker: \gapnote{<kind>}{<status>}. Kind is the policy
% classification (clarification, local-correction, scope-restriction,
% unfaithful, false-source, open-gap); status is open, wip (elimination
% underway), resolved, or historical. Machine-read by the site index;
% prints a verdict line.
\newcommand{\gapnote}[2]{\par\noindent\textbf{Verdict:} #1 (#2).\par}


\title{Length-Dependent Scalars Do Not Determine Copy Weights in the
Proportional BNT Theorem}
\author{}
\date{2026-07-26}

\begin{document}
\maketitle
\gapnote{clarification}{resolved}

\section{The source statement}

Cirac--P\'erez-Garc\'ia--Schuch--Verstraete state that proportional
matrix-product-vector families in canonical form have bijectively matched
bases of normal tensors
\cite[Theorem~2.10 and Appendix~A, lines~1167--1183]
{Cirac2016MPDO_arXiv}. For each matched pair, there are a phase $\zeta_k$ and
an invertible matrix $X_k$ such that
\begin{align}
    Q_k^i
        &=\zeta_kX_kP_{\beta(k)}^iX_k^{-1}.
    \label{eq:cpsv16_proportional_length_scalar_gap_sector_match}
\end{align}
The theorem does not claim that the two total tensors are conjugate, even up
to one common phase. Copy-weight recovery and the block-diagonal global gauge
belong to the equal-MPV Corollary~2.11 of
Ref.~\cite{Cirac2016MPDO_arXiv}, not to the proportional theorem.

The formal proportional theorem assumes eventual nonzero projective
proportionality. Thus, for every sufficiently large $N$, one scalar
$c_N\neq 0$ satisfies
\begin{align}
    V^{(N)}(P_{\mathrm{tot}})
        &=c_NV^{(N)}(Q_{\mathrm{tot}}).
    \label{eq:cpsv16_proportional_length_scalar_gap_total_proportionality}
\end{align}
The source assumption at every positive length is a special case.\footnote{The
formal proportional theorem is
\leanid{MPSTensor.fundamentalTheorem_proportional_canonicalForm} in
\doclink{TNLean/MPS/FundamentalTheorem/SectorBNT/FundamentalCoord}.}

\section{The coefficient identity retaining the global scalar}

For each sector, define
\begin{align}
    c_N^{(j)}(P)
        &=\sum_q\mu_{j,q}^N,
    \label{eq:cpsv16_proportional_length_scalar_gap_p_coefficient}\\
    c_N^{(k)}(Q)
        &=\sum_q\nu_{k,q}^N.
    \label{eq:cpsv16_proportional_length_scalar_gap_q_coefficient}
\end{align}
Substitute
(\ref{eq:cpsv16_proportional_length_scalar_gap_sector_match}) into
(\ref{eq:cpsv16_proportional_length_scalar_gap_total_proportionality}),
reindex by $\beta$, and apply eventual linear independence of the $P$-basis.
For every sector $k$ and all sufficiently large $N$, this gives
\begin{align}
    c_N^{(\beta(k))}(P)
        &=c_N\zeta_k^Nc_N^{(k)}(Q).
    \label{eq:cpsv16_proportional_length_scalar_gap_threaded_coefficient_identity}
\end{align}
The same $c_N$ appears in every sector because it relates the two total
vectors.\footnote{The formal coefficient identity retaining the global scalar
is
\leanid{MPSTensor.coeff_identity_via_matched_mpv_phase_proportional} in
\doclink{TNLean/MPS/FundamentalTheorem/SectorBNT/CoeffIdentity}.}

In the equal-MPV case, $c_N=1$, and the identity becomes
\begin{align}
    \sum_q\mu_{\beta(k),q}^N
        &=\sum_q(\zeta_k\nu_{k,q})^N.
    \label{eq:cpsv16_proportional_length_scalar_gap_scalar_free_power_sums}
\end{align}
This scalar-free identity determines the copy-weight multiset. For a general
proportional family,
(\ref{eq:cpsv16_proportional_length_scalar_gap_threaded_coefficient_identity})
does not imply
(\ref{eq:cpsv16_proportional_length_scalar_gap_scalar_free_power_sums}).

\section{Counterexample}

Let $C$ be a normalized normal tensor. Consider the one-sector BNT forms
\begin{align}
    P
        &=C\oplus\tfrac12C,
    \label{eq:cpsv16_proportional_length_scalar_gap_counterexample_p}\\
    Q
        &=C\oplus\tfrac13C.
    \label{eq:cpsv16_proportional_length_scalar_gap_counterexample_q}
\end{align}
Both satisfy the canonical weight normalization: every weight has modulus at
most one, and each form has a unit-modulus weight. Their MPV families are
\begin{align}
    V^{(N)}(P)
        &=(1+2^{-N})V^{(N)}(C),
    \label{eq:cpsv16_proportional_length_scalar_gap_counterexample_p_mpv}\\
    V^{(N)}(Q)
        &=(1+3^{-N})V^{(N)}(C).
    \label{eq:cpsv16_proportional_length_scalar_gap_counterexample_q_mpv}
\end{align}
They are nonzero and proportional at every positive length, with
\begin{align}
    c_N
        &=\frac{1+2^{-N}}{1+3^{-N}}.
    \label{eq:cpsv16_proportional_length_scalar_gap_counterexample_scalar}
\end{align}
This sequence is not of the form
\begin{align}
    c_N
        &=\xi^N
    \label{eq:cpsv16_proportional_length_scalar_gap_geometric_phase_form}
\end{align}
for a fixed phase $\xi$.

For a one-dimensional choice of $C$, the total matrices have eigenvalue
multisets
\begin{align}
    \operatorname{spec}(P)
        &=\left\{1,\tfrac12\right\},
    \label{eq:cpsv16_proportional_length_scalar_gap_p_spectrum}\\
    \operatorname{spec}(Q)
        &=\left\{1,\tfrac13\right\}.
    \label{eq:cpsv16_proportional_length_scalar_gap_q_spectrum}
\end{align}
Similarity preserves eigenvalues, while multiplication by a common phase
multiplies every eigenvalue by that phase. The total tensors are therefore not
conjugate, even up to a common phase.

\section{Conclusion}

The prior mathematical assessment required correction; there is no
mathematical gap in the source proportional theorem. That theorem does not
require copy-weight recovery or a global gauge, and the counterexample proves
that this stronger claim is false under its hypotheses. In particular, no
missing theorem should assert that the proportionality scalar between two BNT
canonical forms has the geometric phase form
(\ref{eq:cpsv16_proportional_length_scalar_gap_geometric_phase_form}).

The valid conclusions are:
\begin{itemize}
    \item proportional total MPVs determine a bijection and gauge-phase
          equivalence of the normal sectors;
    \item the coefficient identity retains one arbitrary length-dependent
          global scalar $c_N$;
    \item scalar-free coefficient identities, copy-weight matching, and a
          global block gauge require additional hypotheses;
    \item in the equal-MPV corollary, $c_N=1$, so power-sum recovery and the
          global gauge follow.
\end{itemize}

Conditional implications from scalar-free coefficient identities to a global
gauge remain valid.\footnote{The formal conditional global-gauge construction
is
\leanid{MPSTensor.ft_sector_bnt_proportional_global_gauge_of_coeff_identity}
in \doclink{TNLean/MPS/FundamentalTheorem/SectorBNT/Fundamental}.}
They are additional consequences, not unfinished parts of Theorem~2.10 of
Ref.~\cite{Cirac2016MPDO_arXiv}.

\bibliographystyle{alpha}
\bibliography{references}

\end{document}